\documentclass[11pt, a4paper]{article}

% ─── Paquetes ────────────────────────────────────────────────────────────────
\usepackage[utf8]{inputenc}
\usepackage[T1]{fontenc}
\usepackage[spanish]{babel}
\usepackage[top=1.8cm, bottom=1.8cm, left=2cm, right=2cm]{geometry}
\usepackage{graphicx}
\usepackage{enumitem}
\usepackage{titlesec}
\usepackage{xcolor}
\usepackage{hyperref}
\usepackage{microtype}
\usepackage{parskip}
\usepackage{lmodern}
\usepackage{tabularx}
\usepackage{array}
\usepackage{multicol}
\usepackage{fontawesome5}

% ─── Colores ─────────────────────────────────────────────────────────────────
\definecolor{primary}{HTML}{1A237E}
\definecolor{accent}{HTML}{3F51B5}
\definecolor{darkgray}{HTML}{424242}
\definecolor{lightgray}{HTML}{F5F5F5}
\definecolor{rulecol}{HTML}{C5CAE9}

% ─── Hyperref ────────────────────────────────────────────────────────────────
\hypersetup{
  colorlinks=true,
  urlcolor=accent,
  linkcolor=accent,
  pdftitle={CV — Jose Andres Daza Gallego},
  pdfauthor={Jose Andres Daza Gallego},
}

\pagestyle{empty}

% ─── Secciones ───────────────────────────────────────────────────────────────
\titleformat{\section}
  {\large\bfseries\color{primary}}{}{0em}{}[{\color{rulecol}\titlerule[1pt]}]
\titlespacing{\section}{0pt}{10pt}{5pt}

% ─── Listas compactas ────────────────────────────────────────────────────────
\setlist[itemize]{leftmargin=1.4em, itemsep=1pt, topsep=2pt, parsep=0pt}
\setlist[enumerate]{leftmargin=1.4em, itemsep=1pt, topsep=2pt, parsep=0pt}

% ─── Comando para entrada de experiencia ─────────────────────────────────────
\newcommand{\expentry}[4]{%
  \noindent
  \begin{tabularx}{\textwidth}{@{}X r@{}}
    \textbf{\color{primary}#1} & \textit{\small\color{darkgray}#3} \\
    \textit{\color{accent}#2}  & \textit{\small\color{darkgray}#4} \\
  \end{tabularx}%
}

% ─── Comando para entrada de educación ───────────────────────────────────────
\newcommand{\eduentry}[4]{%
  \noindent
  \begin{tabularx}{\textwidth}{@{}X r@{}}
    \textbf{#1} & \textit{\small\color{darkgray}#3 — #4} \\
    \textit{\color{darkgray}#2} & \\
  \end{tabularx}%
}

% ═════════════════════════════════════════════════════════════════════════════
\begin{document}

% ─── Encabezado ──────────────────────────────────────────────────────────────
\begin{center}
  {\Huge\bfseries\color{primary} Jose Andres Daza Gallego}\\[4pt]
  {\large\color{accent} Estudiante de Ingeniería Informática · Técnico en Programación}\\[8pt]
  \small
  \faPhone\ Disponible bajo solicitud \quad
  \faEnvelope\ \href{mailto:foxy200442@gmail.com}{foxy200442@gmail.com} \quad
  \faGithub\ \href{https://github.com/FoxyYTs}{github.com/FoxyYTs}\\[3pt]
  \faLinkedin\ \href{https://www.linkedin.com/in/jose-andres-daza-gallego/}{linkedin.com/in/jose-andres-daza-gallego} \quad
  \faGlobe\ \href{https://foxyyts.github.io/gitprofile/}{foxyyts.github.io/gitprofile}
\end{center}

\vspace{2pt}
{\color{rulecol}\rule{\textwidth}{1.5pt}}
\vspace{2pt}

% ─── Resumen ─────────────────────────────────────────────────────────────────
\section{Resumen de Calificaciones}

Estudiante de Ingeniería Informática con experiencia práctica en desarrollo de software
multi-plataforma, administración de servidores y gestión de laboratorios. Creador de
\textbf{SGI LAB MANAGER}, sistema integral de gestión de laboratorio en producción desarrollado
con Flutter y Django, y de \textbf{Vigía}, sistema de alerta temprana de incendios
forestales con datos satelitales. Capacidad demostrada para liderar proyectos de principio a fin,
trabajar con equipos multidisciplinarios y aplicar tecnologías modernas en entornos reales.

% ─── Habilidades ─────────────────────────────────────────────────────────────
\section{Habilidades Técnicas}

\begin{tabularx}{\textwidth}{@{} >{\bfseries\color{darkgray}}l X @{}}
Lenguajes            & Python, Java, JavaScript, TypeScript, Dart, C++, C\# \\[2pt]
Front-End            & HTML, CSS, Flutter, D3.js \\[2pt]
Back-End             & Django, Celery, Node.js, PHP, pytest \\[2pt]
Bases de Datos       & MySQL, PostgreSQL, PostGIS, MongoDB, Redis \\[2pt]
DevOps / Infra       & Git, Docker, Kubernetes, Nginx, Gunicorn, Cloudflare \\[2pt]
3D / VR              & Unity, Blender, OpenXR, Meta XR SDK \\[2pt]
IA / Datos           & scikit-learn, Groq (Llama 3.1), Google Cloud TTS \\[2pt]
Metodologías         & Scrum \\[2pt]
Sistemas Operativos  & Linux (Ubuntu, Kali, Arch), Bash, Windows \\
\end{tabularx}

\medskip
\noindent\textbf{\color{darkgray}Habilidades Blandas:}
Liderazgo, comunicación efectiva, resolución de problemas, trabajo en equipo.

% ─── Experiencia ─────────────────────────────────────────────────────────────
\section{Experiencia}

\expentry{SGI LAB MANAGER}{Creador y Desarrollador Principal}{Febrero 2024}{Presente}

\begin{itemize}
  \item Diseñé, desarrollé y lancé la primera versión del sistema integral de gestión del
        laboratorio del Politécnico Colombiano Jaime Isaza Cadavid.
  \item Arquitectura multi-plataforma con \textbf{Flutter/Dart} (Android, Web, Windows, Linux)
        y backend \textbf{Django REST Framework} sobre PostgreSQL 15.
  \item DevOps completo: Docker Compose, Nginx, Gunicorn, Cloudflare Zero Trust Tunnel,
        respaldos cifrados automáticos (GPG AES-256) y monitoreo con alertas Telegram.
  \item Módulos: inventario GHS con extracción de FDS por IA (Groq/Llama 3.1), préstamos,
        horario semanal, 6 tipos de informes PDF, formularios públicos por QR e impresiones 3D
        (encargos de varias piezas, abonos y seguimiento por QR).
  \item Más de 450 tests Django, 60+ pytest e2e y 190+ tests Flutter. v1.2.0 (nov 2026).
        Gestión completa del ciclo de vida: planificación, arquitectura, desarrollo,
        despliegue y mantenimiento.
\end{itemize}

\vspace{6pt}
\expentry{Politécnico Colombiano Jaime Isaza Cadavid}{Monitor de Laboratorio}{Febrero 2024}{Presente}

\begin{itemize}
  \item Integro el equipo de monitoría del Laboratorio Integrado junto a otra monitora;
        anteriormente el soporte recaía exclusivamente en las dos laboratoristas titulares.
  \item Brindo soporte técnico diario a estudiantes y docentes, y gestiono el mantenimiento
        preventivo y correctivo de equipos y software especializado.
  \item Las necesidades operativas detectadas en este rol dieron origen al desarrollo de
        \textbf{SGI LAB MANAGER}, sistema hoy en producción activa.
\end{itemize}

\vspace{6pt}
\expentry{Politécnico Colombiano Jaime Isaza Cadavid}{Administrador de Servidores}{Enero 2024}{Presente}

\begin{itemize}
  \item Configuré y administro el servidor institucional que aloja el entorno de
        producción de \textbf{SGI LAB MANAGER}: Docker Compose, Nginx, Gunicorn y
        Cloudflare Zero Trust Tunnel.
  \item Mantengo un homelab personal que replica la infraestructura de producción,
        utilizado como entorno de desarrollo y pruebas antes de desplegar cambios
        en el servidor institucional.
\end{itemize}

\vspace{6pt}
\expentry{Semillero de Ambientes Inmersivos de Aprendizaje — Politécnico Colombiano Jaime Isaza Cadavid}{Investigador y Desarrollador VR}{2023}{Agosto 2026}

\textit{\small\color{darkgray}Proyecto finalizado por falta de apoyo del área de investigación institucional.}

\begin{itemize}
  \item Miembro co-fundador del semillero desde su formación en octubre de 2023.
  \item Desarrollé entornos 3D inmersivos en Unity y VRChat con Meta Quest 2,
        incluyendo el modelo navegable del Bloque R27 y el campus virtual del Politécnico.
  \item Los activos 3D producidos en el semillero son la base del actual proyecto
        Laboratorio Virtual VR.
\end{itemize}

% ─── Proyectos ───────────────────────────────────────────────────────────────
\section{Proyectos Destacados}

\noindent\textbf{\color{primary}SGI LAB MANAGER} \hfill
\href{https://github.com/FoxyYTs/SGI_LAB_MANAGER_FRONTEND}{\small\faGithub\ FoxyYTs/SGI\_LAB\_MANAGER\_FRONTEND}

Plataforma integral de gestión de laboratorio multi-plataforma (Android, Web, Windows, Linux)
con Flutter/Dart y Django REST Framework. Infraestructura propia: Docker Compose, Nginx,
Gunicorn, PostgreSQL 15, Cloudflare Zero Trust Tunnel y respaldos automáticos cifrados
(GPG AES-256) con monitoreo por Telegram. Módulos: inventario GHS con extracción de FDS
por IA (Groq/Llama 3.1), préstamos, horario semanal, 6 tipos de informes PDF y formularios
públicos por QR e impresiones 3D. Más de 450 tests Django · 60+ pytest e2e · 190+ tests Flutter.
v1.2.0 (nov 2026).

\vspace{6pt}
\noindent\textbf{\color{primary}Vigía — Alertas de Incendios Forestales} \hfill
\href{https://github.com/FoxyYTs/Vigia}{\small\faGithub\ FoxyYTs/Vigia} \quad
\href{https://vigia.foxyyts.qzz.io}{\small\faGlobe\ Sitio de producción}

Sistema de alerta temprana de incendios forestales en Colombia, desarrollado de forma
individual como Proyecto Integrador. Ingesta cada 3 horas focos de NASA FIRMS e INPE
(más de 880 mil focos) y cruza históricos de UNGRD e IDEAM; las alertas se agrupan y
llegan por Telegram. Modelo de zonas recurrentes con scikit-learn sobre 1.122 municipios.
Django REST, PostGIS y Celery, con app Flutter web y Android. En producción con Docker y
Cloudflare Tunnel, CI en GitHub Actions y más de 240 tests.

\vspace{6pt}
\noindent\textbf{\color{primary}Pensum Interactivo — PCJIC Rionegro} \hfill
\href{https://github.com/FoxyYTs/pensum-interactivo}{\small\faGithub\ FoxyYTs/pensum-interactivo} \quad
\href{https://foxyyts.github.io/pensum-interactivo}{\small\faGlobe\ Sitio de producción}

Visualizador interactivo de pensum académico para las 8 carreras de la sede Rionegro del
Politécnico Colombiano Jaime Isaza Cadavid. Grafo dirigido con D3.js: asignaturas,
prerrequisitos, créditos y porcentaje de avance con persistencia en \texttt{localStorage}.
PWA desplegada en GitHub Pages con CI/CD automatizado via GitHub Actions. Fork adaptado
de un proyecto chileno de código abierto, con migración de terminología, datos completos
del PCJIC y nuevas funcionalidades.

\vspace{6pt}
\noindent\textbf{\color{primary}Laboratorio Virtual VR} \hfill
\textit{\small\color{darkgray}En desarrollo — Junio 2026}

Sucesor del Semillero AIA. Plataforma de laboratorio
virtual para los Meta Quest 2 del Politécnico, que reutiliza los espacios del Bloque R27
digitalizados durante el semillero como escenario para prácticas interactivas. 4 módulos
(Física, Ondas y Óptica, Química, Electricidad y Magnetismo) con 12 prácticas basadas
en las guías universitarias más una práctica VR-exclusiva por módulo. Unity 6.3 LTS,
C\#, OpenXR, Meta XR SDK. Triple modalidad: Meta Quest 2 standalone, PC VR y escritorio.

% ─── Educación ───────────────────────────────────────────────────────────────
\section{Educación}

\eduentry{Ingeniería Informática \textmd{--- en curso, 9.º semestre}}{Politécnico Colombiano Jaime Isaza Cadavid}{Agosto 2022}{Graduación prevista 2027}

\vspace{4pt}
\eduentry{Técnico en Programación}{I.E. INEM José Félix de Restrepo}{Febrero 2015}{Diciembre 2020}

% ─── Educación adicional ─────────────────────────────────────────────────────
\section{Educación Adicional y Certificaciones}

\begin{itemize}[itemsep=1pt]
  \item \textbf{Git y GitHub} — Platzi (Dic 2025)
  \item \textbf{Fundamentos de Ingeniería de Software} — Platzi (Dic 2025)
  \item \textbf{MongoDB for SQL Professionals} — MongoDB Academy (Nov 2025)
  \item \textbf{Linux Essentials} — Linux Professional Institute / Cisco (Nov 2025)
  \item \textbf{Django} — Platzi (Ago 2025)
  \item \textbf{Python} — Platzi (Ago 2025)
  \item \textbf{Certified Kubernetes Administrator (CKA) with Practice Tests} — Udemy / KodeKloud (Abr 2025)
  \item \textbf{PL DEVOPS} — SoftServe Academy (Jun 2025)
  \item \textbf{Java SE Orientado a Objetos} — Platzi (Sep 2023)
  \item \textbf{Programación Orientada a Objetos} — Platzi (Abr 2023)
  \item \textbf{Desarrollo Personal G8 - ONE} — Alura y Oracle (Ene 2025)
  \item \textbf{Creación de Videojuegos} — Platzi (Ene 2021)
\end{itemize}

% ─── Idiomas ─────────────────────────────────────────────────────────────────
\section{Idiomas}

\noindent\textbf{Español} — Nativo \hfill \textbf{Inglés} — Intermedio (competencia lectora y técnica)

% ─── Referencias ─────────────────────────────────────────────────────────────
\section{Referencias}

\begin{itemize}[leftmargin=0pt, label={}, itemsep=6pt]
  \item \textbf{Luz Delly Jaramillo Carmona} \textit{— Supervisora de Laboratorio, Politécnico Colombiano JIC}\\
        \href{mailto:luzjaramillo@elpoli.edu.co}{luzjaramillo@elpoli.edu.co}
  \item \textbf{Ángela Viviana Martínez Betancur} \textit{— Supervisora de Laboratorio, Politécnico Colombiano JIC}\\
        \href{mailto:angelamartinez@elpoli.edu.co}{angelamartinez@elpoli.edu.co}
  \item \textbf{Ariel Yesid Aristizábal Álzate} \textit{— Supervisor de Laboratorio, Politécnico Colombiano JIC}\\
        \href{mailto:arielaristizabal@elpoli.edu.co}{arielaristizabal@elpoli.edu.co}
\end{itemize}
\small\textit{Teléfonos disponibles bajo solicitud.}\normalsize

% ─── Hobbies ─────────────────────────────────────────────────────────────────
\section{Intereses y Hobbies}

\noindent
\textbf{Electrónica e IoT}~· \textbf{Mecánica de Motocicletas}~· \mbox{\textbf{Desarrollo de Videojuegos}}~· \mbox{\textbf{Proyectos DIY}}~· \textbf{Ciberseguridad}

\end{document}
